\documentclass[11pt]{article}
\usepackage[notes]{calc}
\calcmodule{9}{1}{Basic Ideas}
\calcxdoc{6.1}
\calcxdoc{8.1}
\calcxdoc{8.2}
\calcxdoc{8.5}

\begin{document}

\section{Basic Ideas}

\begin{guide}
In Section~6.1 you were given a velocity and asked for a position: you knew the \emph{derivative} of an unknown function and wanted the function. That is already a differential equation. In general a differential equation is an equation whose unknown is a \emph{function}, and which links the function to its own derivatives. Such equations describe how things change: a cooling cup of coffee, a harvested forest, a draining tank, a falling stone. This section gives the vocabulary (order, linear, solution, general solution, initial value problem), shows how to \emph{check} a proposed solution by substitution, and solves the simplest equations, those of the form $y'=f(t)$, by integration. We prove what the integration method really gives and why an equation of order $n$ has $n$ arbitrary constants. We also watch the domain of a solution: it is often smaller than the domain of the formula. Methods for harder equations come in later sections of this chapter.
\end{guide}

\textbf{By the end of this section you should be able to:}
\begin{itemize}
  \item state the order of a differential equation and decide whether it is linear;
  \item verify by substitution that a function (or a family of functions) is a solution, including the initial conditions;
  \item find the general solution and the solution of an initial value problem for $y^{(n)}=f(t)$ by integration;
  \item determine the domain of a solution from the equation, the initial condition and the context;
  \item set up and check simple models: motion under gravity, a harvested resource, a draining tank;
  \item prove each theorem of this section.
\end{itemize}

% ---------------------------------------------------------------------
\subsection*{What is a differential equation?}

\begin{defn}{Differential equation, order, linear}{de}
A \emph{differential equation} is an equation that contains an unknown function $y$ of one independent variable (usually $t$ or $x$) together with one or more of its derivatives. Its \emph{order} is the highest order of a derivative that appears. It is \emph{linear} if it can be written
\[
a_n(t)\,y^{(n)}+\dots+a_1(t)\,y'+a_0(t)\,y=g(t),
\]
where $a_0,\dots,a_n$ and $g$ are given functions of the independent variable: $y$ and its derivatives appear only to the first power, are not multiplied together, and are not inside another function. Otherwise the equation is \emph{nonlinear}.
\end{defn}

The coefficients $a_j(t)$ may be as complicated as you like; only the dependence on $y$ matters. The unknown is not a number, as in $x^2=4$, but a function.

\begin{exmp}{Order and linearity}{e1}
State the order of each equation and whether it is linear.
\begin{center}
(a) $\dfrac{dy}{dx}+4y=\cos x$ \quad (b) $y''+16y=0$ \quad (c) $y'(t)=0.1\,y\,(100-y)$\\[4pt]
(d) $t^2y'''-2t\,y'+e^ty=\ln t$ \quad (e) $y''+\sin y=0$ \quad (f) $y\,y'=t$
\end{center}

\textbf{Solution.} Find the highest derivative, then look at how $y$ and its derivatives enter.
(a) Highest derivative $y'$: order $1$. The terms $y'$ and $4y$ are of the first power and the right side $\cos x$ depends only on $x$: \emph{linear}.
(b) Highest derivative $y''$: order $2$; every term has the form (function of $t$)$\times y$ or (function of $t$)$\times y''$: \emph{linear}.
(c) Order $1$. The right side is $10y-0.1y^2$, which contains $y^2$: \emph{nonlinear}.
(d) The highest derivative is $y'''$: order $3$. The coefficients $t^2$, $-2t$ and $e^t$ depend on $t$ only, and the right side $\ln t$ too: \emph{linear} (variable coefficients are allowed).
(e) Order $2$. The unknown appears inside the sine, $\sin y$: \emph{nonlinear}.
(f) Order $1$. The product $y\,y'$ multiplies the unknown by its derivative: \emph{nonlinear}.
\end{exmp}

\begin{kc}{Order and linearity}{kcord}
State the order of each equation and whether it is linear: (i) $y''+t\,y'=e^{-t}$; (ii) $y'=y+t^2$; (iii) $y'''=y\,y'$.
\end{kc}

A \emph{solution} should make the equation a true statement at every point of some interval. Because a derivative needs a function that is defined on both sides of a point, we always work on an interval.

\begin{defn}{Solution, general solution, initial value problem}{sol}
A \emph{solution} of a differential equation of order $n$ on an interval $I$ is a function $y$, $n$ times differentiable on $I$, such that substituting $y,y',\dots,y^{(n)}$ turns the equation into an identity for every $t$ in $I$. The interval $I$ is the \emph{domain of the solution}.

The \emph{general solution} is the family of all solutions, written with $n$ arbitrary constants. \emph{Initial conditions} specify the values of $y,y',\dots,y^{(n-1)}$ at one point $t_0$. A differential equation of order $n$ together with $n$ initial conditions is an \emph{initial value problem} (IVP). A \emph{solution of the IVP} is a function that solves the equation on an interval $I$ containing $t_0$ and meets the initial conditions. For $n=1$ the IVP is
\[
y'(t)=F(t,y),\qquad y(t_0)=y_0,
\]
and the solution curve is a graph through the point $(t_0,y_0)$.
\end{defn}

\subsection*{Verifying a solution}

To \emph{check} a candidate you do not need any method of solving: you differentiate and substitute.

\begin{proc}{Verifying a solution by substitution}{ver}
Given a differential equation and a function $y$ (or a family of functions):
\begin{enumerate}[nosep]
  \item Compute every derivative that occurs in the equation.
  \item Substitute into the left side and into the right side; simplify each until they are equal for every $t$ in the interval.
  \item For an IVP, also check each initial condition.
  \item State the interval on which the check is valid.
\end{enumerate}
\end{proc}

\begin{exmp}{A family of solutions and an initial value problem}{e2}
(a) Show that $y(t)=Ce^{-3t}$ solves $y'=-3y$ for every constant $C$. (b) Solve the IVP $y'=-3y$, $y(0)=2$.

\textbf{Solution.} (a) By the Chain Rule, $y'(t)=C\cdot(-3)e^{-3t}=-3Ce^{-3t}$. The right side is $-3y(t)=-3Ce^{-3t}$. The two sides agree for every real $t$ and every $C$, so $y=Ce^{-3t}$ is a family of solutions on $(-\infty,\infty)$.
%@ e2-fam | simplify(diff(c*exp(-3*t),t) - (-3)*(c*exp(-3*t))) == 0 | $-3y(t)=-3Ce^{-3t}$. The two sides agree for every real $t$ and every $C$,
(b) The initial condition fixes $C$: $y(0)=Ce^{0}=C=2$. The solution is $y(t)=2e^{-3t}$. Check: $y'(t)=-6e^{-3t}=-3\cdot2e^{-3t}=-3y(t)$ and $y(0)=2$. The figure shows several members of the family; the IVP singles out the one through $(0,2)$.
%@ e2-ivp | simplify(diff(2*exp(-3*t),t) - (-3)*2*exp(-3*t)) == 0 and (2*exp(-3*t)).subs(t,0) == 2 | and $y(0)=2$. The figure shows several members of the family; the IVP
\begin{center}
\begin{tikzpicture}[x=3.2cm,y=0.55cm]
  \draw[->] (-0.1,0) -- (1.75,0) node[right] {$t$};
  \draw[->] (0,-2.7) -- (0,4.6) node[above] {$y$};
  \foreach \C in {-2,-1,1,4}{\draw[thin,gray,domain=0:1.6,samples=40,variable=\t] plot ({\t},{\C*exp(-3*\t)});}
  \draw[very thick,blue!70!black,domain=0:1.6,samples=40,variable=\t] plot ({\t},{2*exp(-3*\t)});
  \fill (0,2) circle (1.6pt) node[left,font=\small] {$(0,2)$};
%@ g-e2-p | 2*exp(-3*0) == 2 | circle (1.6pt) node[left,font=\small] {$(0,2)$};
  \fill (0,4) circle (1.2pt) node[left,font=\small] {$C=4$};
%@ g-e2-c4 | 4*exp(-3*0) == 4 | circle (1.2pt) node[left,font=\small] {$C=4$};
  \fill (0,-2) circle (1.2pt) node[left,font=\small] {$C=-2$};
%@ g-e2-cm | -2*exp(-3*0) == -2 | circle (1.2pt) node[left,font=\small] {$C=-2$};
\end{tikzpicture}
\par\small\emph{Figure.} Solutions $y=Ce^{-3t}$ for $C=-2,-1,1,2,4$; the blue curve is the solution of the IVP.
\end{center}
\end{exmp}

\textbf{Linear equations have a structure.} In Example~\ref{ex:e2}, if $y_1=e^{-3t}$ is a solution then so is every constant multiple. This is no accident, and it is the reason the general solution of a linear equation can be built from a few special solutions.

\begin{thm}{Sums and multiples of solutions of a linear equation}{lin}
Let $L[y]=y^{(n)}+p_{n-1}(t)\,y^{(n-1)}+\dots+p_1(t)\,y'+p_0(t)\,y$, where the $p_j$ are given functions, and let $y_1,y_2$ be $n$ times differentiable on an interval $I$. For all constants $c_1,c_2$:
\begin{enumerate}[label=(\alph*),nosep]
  \item $L[c_1y_1+c_2y_2]=c_1L[y_1]+c_2L[y_2]$;
  \item if $L[y_1]=L[y_2]=0$, then $L[c_1y_1+c_2y_2]=0$;
  \item if $L[y_p]=q$ and $L[y_h]=0$, then $L[y_p+y_h]=q$;
  \item if $L[y_1]=L[y_2]=q$, then $L[y_1-y_2]=0$.
\end{enumerate}
\end{thm}
\begin{proof}
(a) By the sum and constant-multiple rules for derivatives, $(c_1y_1+c_2y_2)^{(j)}=c_1y_1^{(j)}+c_2y_2^{(j)}$ for each $j=0,1,\dots,n$. Multiply the $j$th identity by $p_j$ (with $p_n=1$) and add over $j$; the terms regroup as $c_1L[y_1]+c_2L[y_2]$.
(b) Put $L[y_1]=L[y_2]=0$ in (a). (c) Apply (a) with $c_1=c_2=1$: $L[y_p+y_h]=q+0=q$. (d) Apply (a) with $c_1=1$, $c_2=-1$: $L[y_1-y_2]=q-q=0$.
\end{proof}

The equation $L[y]=0$ is called \emph{homogeneous}; part (c) says a solution of $L[y]=q$ plus any solution of the homogeneous equation is again a solution. The proof used the \emph{form} of $L$, so nothing like it works for a nonlinear equation (Pause and think~\ref{think:thnl} below).

\begin{exmp}{A second-order equation with two constants}{e3}
(a) Show that $y=C_1\cos4t+C_2\sin4t$ is a solution of $y''+16y=0$ for all constants $C_1,C_2$. (b) Solve the IVP $y''+16y=0$, $y(0)=2$, $y'(0)=-4$.

\textbf{Solution.} (a) Check the two special functions, then use Theorem~\ref{thm:lin}(b) with $L[y]=y''+16y$. For $y_1=\cos4t$: $y_1'=-4\sin4t$, $y_1''=-16\cos4t$, so $y_1''+16y_1=-16\cos4t+16\cos4t=0$. For $y_2=\sin4t$: $y_2'=4\cos4t$, $y_2''=-16\sin4t$, and $y_2''+16y_2=0$.
%@ e3-a | simplify(diff(cos(4*t),t,2)+16*cos(4*t)) == 0 and simplify(diff(sin(4*t),t,2)+16*sin(4*t)) == 0 | $y_2''=-16\sin4t$, and $y_2''+16y_2=0$.
By the theorem, $C_1y_1+C_2y_2$ is a solution for all $C_1,C_2$, on $(-\infty,\infty)$.
(b) Differentiate the family: $y'(t)=-4C_1\sin4t+4C_2\cos4t$. At $t=0$: $y(0)=C_1\cos0+C_2\sin0=C_1=2$ and $y'(0)=4C_2=-4$, so $C_2=-1$. The solution is $y(t)=2\cos4t-\sin4t$. Check: $y''=-32\cos4t+16\sin4t=-16y$, so $y''+16y=0$; $y(0)=2$; $y'(0)=-8\sin0-4\cos0=-4$.
%@ e3-b | simplify(diff(2*cos(4*t)-sin(4*t),t,2)+16*(2*cos(4*t)-sin(4*t))) == 0 and (2*cos(4*t)-sin(4*t)).subs(t,0) == 2 and diff(2*cos(4*t)-sin(4*t),t).subs(t,0) == -4 | so $y''+16y=0$; $y(0)=2$; $y'(0)=-8\sin0-4\cos0=-4$.
\end{exmp}

An equation of order $n$ has $n$ constants, so it needs $n$ initial conditions: here two, a position and a velocity. Why that count is correct is explained below for the simplest equations.

% ---------------------------------------------------------------------
\subsection*{Solving $y^{(n)}=f(t)$ by integration}

When the right side does not involve $y$, the equation $y'=f(t)$ just asks for an antiderivative, as in Section~6.1.

\begin{thm}{The equation $y'=f(t)$}{anti}
Let $f$ be continuous on an interval $I$ and let $t_0\in I$.
\begin{enumerate}[label=(\alph*),nosep]
  \item If $F$ is one antiderivative of $f$ on $I$, then the general solution of $y'=f(t)$ on $I$ is $y=F(t)+C$: every function of this form is a solution, and every solution has this form.
  \item For every number $y_0$, the IVP $y'=f(t)$, $y(t_0)=y_0$ has exactly one solution on $I$, namely
  \[
  y(t)=y_0+\int_{t_0}^tf(x)\,dx .
  \]
\end{enumerate}
\end{thm}
\begin{proof}
(a) $(F+C)'=f$, so each $F+C$ is a solution. Conversely let $y$ be a solution on $I$. Then $(y-F)'=f-f=0$ on $I$. If $t_1<t_2$ are in $I$, the Mean Value Theorem (Calculus~I) gives $c$ in $(t_1,t_2)$ with $(y-F)(t_2)-(y-F)(t_1)=(y-F)'(c)\,(t_2-t_1)=0$. So $y-F$ takes the same value at any two points of $I$, that is, it is a constant $C$.
(b) Let $Y(t)=y_0+\int_{t_0}^tf(x)\,dx$. By the Fundamental Theorem of Calculus (Part~1) $Y'=f$ on $I$, and $Y(t_0)=y_0+0=y_0$, so $Y$ solves the IVP. If $y$ is any solution of the IVP, then $y=Y+C$ by (a) (with $F=Y$), and $y(t_0)=y_0=Y(t_0)+C$ gives $C=0$, so $y=Y$.
\end{proof}

This is Theorem~\ref{6.1-thm:pos} again, with $t_0$ in place of $0$: the position is the initial position plus the integral of the velocity. There are two ways to find the solution of an IVP, the antiderivative method (find $F+C$, then use the condition to find $C$) and the integral method of part (b); they always agree.

Higher-order equations $y^{(n)}=f(t)$ are solved by integrating $n$ times, and each integration brings one constant.

\begin{thm}{The equation $y^{(n)}=f(t)$}{nth}
Let $f$ be continuous on an interval $I$, $t_0\in I$, and let $b_0,b_1,\dots,b_{n-1}$ be any numbers. Then the IVP
\[
y^{(n)}=f(t),\qquad y(t_0)=b_0,\ y'(t_0)=b_1,\ \dots,\ y^{(n-1)}(t_0)=b_{n-1}
\]
has exactly one solution on $I$. In particular the general solution has exactly $n$ free constants, which can be taken to be $b_0,\dots,b_{n-1}$.
\end{thm}
\begin{proof}
We use induction on $n$. For $n=1$ this is Theorem~\ref{thm:anti}(b). Let $n\ge2$ and assume the statement for $n-1$. A function $y$ solves the IVP of order $n$ exactly when $w=y'$ solves $w^{(n-1)}=f$ with $w(t_0)=b_1,\dots,w^{(n-2)}(t_0)=b_{n-1}$, and $y(t_0)=b_0$. By the induction hypothesis there is exactly one such $w$. By Theorem~\ref{thm:anti}(b) applied to $y'=w$ (the function $w$ is continuous, being differentiable) there is exactly one $y$ with $y'=w$ and $y(t_0)=b_0$. So there is exactly one solution.
\end{proof}

\begin{think}{The constant on a broken domain}{thgap}
The function $F(t)=-\dfrac1t$ is an antiderivative of $f(t)=\dfrac1{t^2}$ on $(-\infty,0)$ and on $(0,\infty)$. Is $y=-\dfrac1t+C_1$ for $t<0$ and $y=-\dfrac1t+C_2$ for $t>0$, with $C_1\ne C_2$, a solution of $y'=\dfrac1{t^2}$ on $t\ne0$? Why does Theorem~\ref{thm:anti} demand an interval?
\end{think}

\begin{exmp}{General solutions by integration}{e4}
Find the general solution: (a) $y'(t)=6\sin2t+8e^{-4t}$; (b) $y''(t)=24t^2-3\cos3t+4$.

\textbf{Solution.} (a) By Theorem~\ref{thm:anti}, the solutions are the antiderivatives. With $u=2t$ and $u=-4t$ (Theorem~\ref{8.1-thm:subst}): $\int6\sin2t\,dt=-3\cos2t$ and $\int8e^{-4t}\,dt=-2e^{-4t}$. So $y(t)=-3\cos2t-2e^{-4t}+C$. Check: $y'=6\sin2t+8e^{-4t}$.
%@ e4-a | ANT(-3*cos(2*x)-2*exp(-4*x), 6*sin(2*x)+8*exp(-4*x), [Q(1,2),Q(3,2),Q(-1,3)]) == 0 | So $y(t)=-3\cos2t-2e^{-4t}+C$. Check: $y'=6\sin2t+8e^{-4t}$.
(b) Integrate once: $y'(t)=\int(24t^2-3\cos3t+4)\,dt=8t^3-\sin3t+4t+C_1$. Integrate again: $y(t)=2t^4+\frac13\cos3t+2t^2+C_1t+C_2$, since $\int(-\sin3t)\,dt=\frac13\cos3t$. Check: $y'=8t^3-\sin3t+4t+C_1$ and $y''=24t^2-3\cos3t+4$. Two integrations, two constants; by Theorem~\ref{thm:nth} these are all the solutions.
%@ e4-b | ANT(2*x**4+cos(3*x)/3+2*x**2, 8*x**3-sin(3*x)+4*x, [Q(1,2),Q(3,2)]) == 0 and ANT(8*x**3-sin(3*x)+4*x, 24*x**2-3*cos(3*x)+4, [Q(1,2),Q(3,2)]) == 0 | and $y''=24t^2-3\cos3t+4$. Two integrations, two constants; by Theorem~\ref{thm:nth} these are all the solutions.
\end{exmp}

\begin{exmp}{Initial value problems, by two methods}{e5}
Solve: (a) $y'(t)=12e^{-t/4}$, $y(0)=7$; (b) $y'(t)=te^{t}$, $y(0)=-1$.

\textbf{Solution.} (a) \emph{Antiderivative method.} $y(t)=\int12e^{-t/4}\,dt=-48e^{-t/4}+C$ (check: $\frac{d}{dt}(-48e^{-t/4})=12e^{-t/4}$). Then $y(0)=-48+C=7$, so $C=55$ and $y(t)=55-48e^{-t/4}$. Check: $y(0)=7$. As $t\to\infty$, $e^{-t/4}\to0$ and $y\to55$: the solution rises toward the value $55$ but never reaches it.
%@ e5-a | ANT(-48*exp(-x/4), 12*exp(-x/4), [Q(1,2),Q(3,1)]) == 0 and 55-48 == 7 and limit(55-48*exp(-x/4), x, oo) == 55 | $t\to\infty$, $e^{-t/4}\to0$ and $y\to55$: the solution rises toward the value $55$ but never reaches it.
(b) \emph{Integral method} (Theorem~\ref{thm:anti}(b) with $t_0=0$): $y(t)=-1+\int_0^txe^x\,dx$. Integration by parts (Theorem~\ref{8.2-thm:parts}) with $u=x$, $dv=e^x\,dx$ gives $\int xe^x\,dx=xe^x-e^x$. So
\[
y(t)=-1+\Big[xe^x-e^x\Big]_0^t=-1+(te^t-e^t)-(0-1)=(t-1)e^t .
\]
Check: $y'=e^t+(t-1)e^t=te^t$ and $y(0)=-1$.
%@ e5-b | ANT(x*exp(x)-exp(x), x*exp(x), [Q(1,2),Q(3,2)]) == 0 and simplify(diff((x-1)*exp(x),x)-x*exp(x)) == 0 and expand((x-1)*exp(x) - (-1+(x*exp(x)-exp(x))-(0-1))) == 0 and (0-1)*exp(0) == -1 | Check: $y'=e^t+(t-1)e^t=te^t$ and $y(0)=-1$.
\end{exmp}

\begin{kc}{An initial value problem}{kciv}
Solve $y'(t)=12e^{-t/4}$ with the initial condition $y(0)=-5$ (compare Example~\ref{ex:e5}(a)). What does $y$ approach as $t\to\infty$?
\end{kc}

\subsection*{The domain of a solution}

A formula such as $y=t+3t^{2/3}$ has its own domain, but as a solution of a differential equation it must be differentiable and must satisfy the equation. The domain of the solution is the \emph{interval containing $t_0$} on which this is so.

\begin{exmp}{Domain from the equation}{e6}
Solve each IVP and find the domain of the solution.
(a) $y'(t)=3+\dfrac2{\sqrt[3]t}$, $y(8)=40$. \quad (b) $y'(x)=2\sec2x\tan2x$, $y(\pi/6)=0$.

\textbf{Solution.} (a) Write the equation as $y'=3+2t^{-1/3}$. By the power rule, $y(t)=3t+2\cdot\frac32t^{2/3}+C=3t+3t^{2/3}+C$ (check: $\frac{d}{dt}3t^{2/3}=2t^{-1/3}$). The condition gives $y(8)=24+3\cdot8^{2/3}+C=24+12+C=40$, so $C=4$ and
\[
y(t)=3t+3t^{2/3}+4 .
\]
As a formula the domain is all real $t$. But $y'=3+2t^{-1/3}$ is not defined at $t=0$, so $y$ solves the equation only on an interval that avoids $0$. The initial point $t_0=8$ lies in $(0,\infty)$, so the \emph{domain of the solution is $(0,\infty)$}. The left branch (the same formula for $t<0$) solves the equation for other initial conditions.
%@ e6-a | ANT(3*x+3*x**Q(2,3)+4, 3+2*x**Q(-1,3), [Q(1,2),Q(8,1),Q(27,1)]) == 0 and 24+3*8**Q(2,3)+4 == 40 and 8**Q(2,3) == 4 | domain of the solution is $(0,\infty)$}. The left branch (the same formula for $t<0$) solves the equation for other initial conditions.
(b) With $u=2x$, $du=2\,dx$, and $\frac d{du}\sec u=\sec u\tan u$: $\int2\sec2x\tan2x\,dx=\sec2x+C$. So $y=\sec2x+C$ and $y(\pi/6)=\sec\frac\pi3+C=2+C=0$ gives $C=-2$:
\[
y(x)=\sec2x-2 .
\]
Check: $y'=2\sec2x\tan2x$. Now $\sec2x=1/\cos2x$ is undefined where $\cos2x=0$, that is, at $x=\pm\frac\pi4,\pm\frac{3\pi}4,\dots$. The point $\frac\pi6$ lies between $-\frac\pi4$ and $\frac\pi4$, so the domain of the solution is $\left(-\frac\pi4,\frac\pi4\right)$.
%@ e6-b | ANT(sec(2*x)-2, 2*sec(2*x)*tan(2*x), [Q(1,3),Q(1,2),Q(-1,2)]) == 0 and sec(pi/3)-2 == 0 and cos(2*(pi/4)) == 0 | so the domain of the solution is $\left(-\frac\pi4,\frac\pi4\right)$.

\begin{center}
\begin{tikzpicture}[x=0.26cm,y=0.055cm]
  \draw[->] (-10,0) -- (13.5,0) node[right] {$t$};
  \draw[->] (0,-14) -- (0,60) node[above] {$y$};
  \draw[thin,gray,dashed,domain=-9:-0.01,samples=60,variable=\t] plot ({\t},{3*\t+3*pow(abs(\t),0.666667)+4});
  \draw[very thick,blue!70!black,domain=0.01:12,samples=60,variable=\t] plot ({\t},{3*\t+3*pow(\t,0.666667)+4});
  \fill (8,40) circle (1.8pt) node[below right,font=\small] {$(8,40)$};
%@ g-e6-p | 3*8+3*8**Q(2,3)+4 == 40 | circle (1.8pt) node[below right,font=\small] {$(8,40)$};
  \draw[blue!70!black] (0,4) circle (2pt);
  \node[font=\small,anchor=east] at (-0.5,9) {$(0,4)$};
%@ g-e6-c | 3*0+3*0**Q(2,3)+4 == 4 |   \node[font=\small,anchor=east] at (-0.5,9) {$(0,4)$};
  \fill (-8,-8) circle (1.5pt) node[left,font=\small] {$(-8,-8)$};
%@ g-e6-l | 3*(-8)+3*8**Q(2,3)+4 == -8 | circle (1.5pt) node[left,font=\small] {$(-8,-8)$};
\end{tikzpicture}\hspace{0.6cm}
\begin{tikzpicture}[x=1.0cm,y=0.5cm]
  \draw[->] (-2.5,0) -- (2.8,0) node[right] {$x$};
  \draw[->] (0,-6) -- (0,3.4) node[above] {$y$};
  \draw[thin,dashed,gray] (0.7854,-6) -- (0.7854,3.2);
  \draw[thin,dashed,gray] (-0.7854,-6) -- (-0.7854,3.2);
  \begin{scope}
    \clip (-2.5,-6) rectangle (2.7,3.2);
    \draw[thin,gray,domain=0.82:2.32,samples=80,variable=\x] plot ({\x},{1/cos(2*\x r)-2});
    \draw[thin,gray,domain=-2.32:-0.82,samples=80,variable=\x] plot ({\x},{1/cos(2*\x r)-2});
    \draw[very thick,blue!70!black,domain=-0.75:0.75,samples=60,variable=\x] plot ({\x},{1/cos(2*\x r)-2});
  \end{scope}
  \fill (0.5236,0) circle (1.8pt) node[above left,font=\small] {$\pi/6$};
%@ g-e6-q | sec(2*(pi/6))-2 == 0 | circle (1.8pt) node[above left,font=\small] {$\pi/6$};
  \fill (0,-1) circle (1.5pt) node[below left,font=\small] {$-1$};
%@ g-e6-m | sec(0)-2 == -1 | circle (1.5pt) node[below left,font=\small] {$-1$};
  \node[font=\small,anchor=north west] at (0.7854,0) {$\pi/4$};
%@ g-e6-as | cos(2*(pi/4)) == 0 |   \node[font=\small,anchor=north west] at (0.7854,0) {$\pi/4$};
\end{tikzpicture}
\par\small\emph{Figure.} Left: $y=3t+3t^{2/3}+4$; only the right branch (heavy) solves the IVP, and $y'$ does not exist at $t=0$. Right: $y=\sec2x-2$; only the branch on $(-\frac\pi4,\frac\pi4)$ (heavy) solves the IVP.
\end{center}
\end{exmp}

\begin{kc}{Which branch?}{kcdom}
Solve $y'(t)=3+\dfrac2{\sqrt[3]t}$ with $y(-8)=10$ and state the domain of the solution (compare Example~\ref{ex:e6}(a)).
\end{kc}

Even when the right side of a differential equation is perfectly smooth, the \emph{solution} can fail to exist for all $t$.

\begin{exmp}{A solution that blows up}{e7}
(a) Show that $y=\dfrac1{C-t}$ is a solution of the nonlinear equation $y'=y^2$ on every interval on which $C\ne t$. (b) Solve $y'=y^2$ with $y(0)=1$ and with $y(0)=-1$, and find each domain. (c) Is $y=0$ in the family?

\textbf{Solution.} (a) $y=(C-t)^{-1}$, so by the Chain Rule $y'=-(C-t)^{-2}\cdot(-1)=\dfrac1{(C-t)^2}=y^2$.
%@ e7-a | simplify(diff(1/(c-t),t) - (1/(c-t))**2) == 0 | $y'=-(C-t)^{-2}\cdot(-1)=\dfrac1{(C-t)^2}=y^2$.
(b) If $y(0)=1$: $\frac1C=1$, so $C=1$ and $y=\dfrac1{1-t}$. It is defined for $t\ne1$; the point $t_0=0$ lies in $(-\infty,1)$, so the domain is $(-\infty,1)$, and $y\to\infty$ as $t\to1^-$: the solution \emph{blows up in finite time}. If $y(0)=-1$: $\frac1C=-1$, $C=-1$ and $y=\dfrac1{-1-t}=-\dfrac1{1+t}$, with domain $(-1,\infty)$ (it blows up as $t\to-1^+$, backward in time).
%@ e7-b | limit(1/(1-x), x, 1, '-') == oo and limit(-1/(1+x), x, -1, '+') == -oo and 1/(1-0) == 1 and -1/(1+0) == -1 | and $y=\dfrac1{-1-t}=-\dfrac1{1+t}$, with domain $(-1,\infty)$ (it blows up as $t\to-1^+$, backward in time).
(c) No value of $C$ gives $y=0$, since $\frac1{C-t}\ne0$. But $y=0$ is a solution: $0'=0=0^2$. It is an \emph{isolated} solution not contained in the general solution, which is why we call $y=1/(C-t)$ the general solution only in the sense of ``the family of solutions the method produces''.
%@ e7-c | diff(Integer(0),t) == Integer(0)**2 | $y=0$ is a solution: $0'=0=0^2$. It is an \emph{isolated} solution not contained in the general solution, which is why we call

\begin{center}
\begin{tikzpicture}[x=0.9cm,y=0.5cm]
  \draw[->] (-3.3,0) -- (3.4,0) node[right] {$t$};
  \draw[->] (0,-4.2) -- (0,4.4) node[above] {$y$};
  \draw[thin,dashed,gray] (1,-4.2) -- (1,4.2);
  \draw[thin,dashed,gray] (-1,-4.2) -- (-1,4.2);
  \begin{scope}
    \clip (-3.3,-4.2) rectangle (3.3,4.2);
    \draw[very thick,blue!70!black,domain=-3:0.75,samples=60,variable=\t] plot ({\t},{1/(1-\t)});
    \draw[very thick,red!70!black,domain=-0.75:3,samples=60,variable=\t] plot ({\t},{-1/(1+\t)});
  \end{scope}
  \fill (0,1) circle (1.8pt) node[below right,font=\small] {$(0,1)$};
%@ g-e7-a | 1/(1-0) == 1 | circle (1.8pt) node[below right,font=\small] {$(0,1)$};
  \fill (0,-1) circle (1.8pt) node[above left,font=\small] {$(0,-1)$};
%@ g-e7-b | -1/(1+0) == -1 | circle (1.8pt) node[above left,font=\small] {$(0,-1)$};
  \node[font=\small,anchor=north west] at (1.0,4.2) {$t=1$};
%@ g-e7-t1 | 1-1 == 0 |   \node[font=\small,anchor=north west] at (1.0,4.2) {$t=1$};
  \node[font=\small,anchor=south east] at (-1.0,-4.2) {$t=-1$};
%@ g-e7-t2 | -1+1 == 0 |   \node[font=\small,anchor=south east] at (-1.0,-4.2) {$t=-1$};
\end{tikzpicture}
\par\small\emph{Figure.} $y=1/(1-t)$ on $(-\infty,1)$ and $y=-1/(1+t)$ on $(-1,\infty)$; each stops at its vertical asymptote. The solution $y=0$ is the $t$-axis.
\end{center}
\end{exmp}

\begin{think}{Multiples of solutions}{thnl}
$y_1=\dfrac1{1-t}$ solves the nonlinear equation $y'=y^2$ on $(-\infty,1)$. Is $2y_1$ a solution? For the linear equation $y'=-3y$ the function $2e^{-3t}$ is a multiple of the solution $e^{-3t}$. Which step in the proof of Theorem~\ref{thm:lin} breaks for $y'=y^2$?
\end{think}

\subsection*{Differential equations in action}

Three models show how the domain is settled by the equation and also by the physical context. Here we only \emph{verify} solutions; ways of \emph{finding} them are the subject of later sections of this chapter.

\begin{exmp}{Motion in a gravitational field}{e8}
A stone is launched straight up with speed $v_0>0$ m/s from a height $s_0\ge0$ m at $t=0$. Take $s$ positive upward, and neglect air resistance. By Newton's second law $s''(t)=-g$, with $g=9.8$ m/s$^2$. (a) Find $s(t)$ and the domain of the solution. (b) Find the time and height of the highest point. (c) Which goes higher: the stone with $(v_0,s_0)=(24.5,0)$ or the one with $(9.8,28)$?

\textbf{Solution.} (a) Integrate once: $s'(t)=-9.8t+C_1$. The condition $s'(0)=v_0$ gives $C_1=v_0$. Integrate again: $s(t)=-4.9t^2+v_0t+C_2$, and $s(0)=s_0$ gives $C_2=s_0$. By Theorem~\ref{thm:nth} this is the only solution:
\[
s(t)=-4.9t^2+v_0t+s_0 .
\]
The model is valid only while the stone is in the air, $t\ge0$ and $s(t)\ge0$. The quadratic $-4.9t^2+v_0t+s_0$ has one positive root, $T=\dfrac{v_0+\sqrt{v_0^2+19.6\,s_0}}{9.8}$ (quadratic formula), so the domain is $[0,T]$.
%@ e8-a | simplify(diff(-Q(49,10)*t**2+a*t+b,t,2) + Q(98,10)) == 0 and (-Q(49,10)*t**2+a*t+b).subs(t,0) == b and diff(-Q(49,10)*t**2+a*t+b,t).subs(t,0) == a | The quadratic $-4.9t^2+v_0t+s_0$ has one positive root, $T=\dfrac{v_0+\sqrt{v_0^2+19.6\,s_0}}{9.8}$ (quadratic formula), so the domain is $[0,T]$.
(b) The highest point is where $s'(t)=-9.8t+v_0=0$, that is $t=\dfrac{v_0}{9.8}$. It is a maximum since $s$ is a downward parabola. The height is
\[
s\!\left(\frac{v_0}{9.8}\right)=-4.9\frac{v_0^2}{9.8^2}+\frac{v_0^2}{9.8}+s_0=\frac{v_0^2}{19.6}+s_0 .
\]
%@ e8-b | simplify(-Q(49,10)*(a/Q(98,10))**2+a**2/Q(98,10)+b - (a**2/Q(196,10)+b)) == 0 | s\!\left(\frac{v_0}{9.8}\right)=-4.9\frac{v_0^2}{9.8^2}+\frac{v_0^2}{9.8}+s_0=\frac{v_0^2}{19.6}+s_0 .
(c) With $(24.5,0)$: peak at $t=\frac{24.5}{9.8}=2.5$ s, height $\frac{24.5^2}{19.6}=30.625$ m. With $(9.8,28)$: peak at $t=1$ s, height $\frac{9.8^2}{19.6}+28=4.9+28=32.9$ m. So the slower stone launched from higher up goes higher (see the figure). They land at $T=5$ s and $T=\frac{9.8+\sqrt{9.8^2+19.6\cdot28}}{9.8}=1+\sqrt{\frac{47}7}\approx3.59$ s.
%@ e8-c | Q(245,10)/Q(98,10) == Q(5,2) and Q(245,10)**2/Q(196,10) == Q(245,8) and Q(98,10)**2/Q(196,10)+28 == Q(329,10) and NUM(1+sqrt(Q(47,7)), (Q(98,10)+sqrt(Q(98,10)**2+Q(196,10)*28))/Q(98,10)) == 0 | $T=\frac{9.8+\sqrt{9.8^2+19.6\cdot28}}{9.8}=1+\sqrt{\frac{47}7}\approx3.59$ s.

\begin{center}
\begin{tikzpicture}[x=1.1cm,y=0.13cm]
  \draw[->] (-0.2,0) -- (5.7,0) node[right] {$t$ (s)};
  \draw[->] (0,-2) -- (0,38) node[above] {$s$ (m)};
  \draw[very thick,blue!70!black,domain=0:5,samples=50,variable=\t] plot ({\t},{-4.9*\t*\t+24.5*\t});
  \draw[very thick,red!70!black,domain=0:3.591,samples=50,variable=\t] plot ({\t},{-4.9*\t*\t+9.8*\t+28});
  \fill (2.5,30.625) circle (1.6pt) node[above,font=\small] {$(2.5,30.6)$};
%@ g-e8-a | -Q(49,10)*Q(5,2)**2+Q(245,10)*Q(5,2) == Q(245,8) | circle (1.6pt) node[above,font=\small] {$(2.5,30.6)$};
  \fill (1,32.9) circle (1.6pt) node[above,font=\small] {$(1,32.9)$};
%@ g-e8-b | -Q(49,10)*1+Q(98,10)*1+28 == Q(329,10) | circle (1.6pt) node[above,font=\small] {$(1,32.9)$};
  \fill (0,28) circle (1.6pt) node[left,font=\small] {$28$};
%@ g-e8-c | -Q(49,10)*0+Q(98,10)*0+28 == 28 | circle (1.6pt) node[left,font=\small] {$28$};
  \fill (5,0) circle (1.6pt) node[below,font=\small] {$5$};
%@ g-e8-d | -Q(49,10)*25+Q(245,10)*5 == 0 | circle (1.6pt) node[below,font=\small] {$5$};
  \fill (3.591,0) circle (1.6pt) node[below,font=\small] {$3.59$};
%@ g-e8-e | abs((1+sqrt(Q(47,7))).evalf(30)-Q(359,100)) < Q(1,100) | circle (1.6pt) node[below,font=\small] {$3.59$};
\end{tikzpicture}
\par\small\emph{Figure.} Position against time (not the path of the stone, which moves along a vertical line). Blue: $(v_0,s_0)=(24.5,0)$. Red: $(9.8,28)$.
\end{center}
\end{exmp}

\begin{exmp}{A harvesting model}{e9}
A resource (timber or fish) with amount $p(t)$ grows naturally at the rate $rp$ while being harvested at the constant rate $H>0$: $p'(t)=rp(t)-H$, $p(0)=p_0$, with $r>0$. (a) Verify that $p(t)=\left(p_0-\dfrac Hr\right)e^{rt}+\dfrac Hr$ solves the IVP. (b) Let $r=0.2$ and $p_0=500$. Describe the solutions for $H=60,100,140$. (c) For which $H$ is $p$ constant, and when does the resource vanish if $H=140$?

\textbf{Solution.} (a) Differentiate: $p'(t)=r\left(p_0-\frac Hr\right)e^{rt}=(rp_0-H)e^{rt}$. Substitute into the right side: $rp-H=r\left(p_0-\frac Hr\right)e^{rt}+H-H=(rp_0-H)e^{rt}$. The sides agree. At $t=0$: $p(0)=p_0-\frac Hr+\frac Hr=p_0$.
%@ e9-a | simplify(diff((a-b/r)*exp(r*t)+b/r,t) - (r*((a-b/r)*exp(r*t)+b/r)-b)) == 0 and simplify(((a-b/r)*exp(r*t)+b/r).subs(t,0) - a) == 0 | At $t=0$: $p(0)=p_0-\frac Hr+\frac Hr=p_0$.
(b) Here $\frac Hr=5H$ and $p(t)=(500-5H)e^{0.2t}+5H$. For $H=60$: $p=200e^{0.2t}+300$, increasing. For $H=100$: $500-5H=0$, so $p(t)=500$ for all $t$. For $H=140$: $p=700-200e^{0.2t}$, decreasing.
%@ e9-b | simplify((500-5*60)*exp(t/5)+5*60 - (200*exp(t/5)+300)) == 0 and 500-5*100 == 0 and simplify((500-5*140)*exp(t/5)+5*140 - (700-200*exp(t/5))) == 0 | For $H=140$: $p=700-200e^{0.2t}$, decreasing.
(c) The coefficient $p_0-\frac Hr$ vanishes exactly when $H=rp_0=0.2\cdot500=100$, the \emph{equilibrium} case (harvest balances growth). For $H=140$, $p=0$ when $200e^{0.2t}=700$, so $e^{0.2t}=3.5$ and $t=5\ln3.5\approx6.26$. The domain of the solution is $[0,5\ln3.5]$: after that the formula gives negative amounts, which have no meaning.
%@ e9-c | 0.2*500 == 100 and 700-200*exp(5*log(Q(7,2))/5) == 0 and abs((5*log(Q(7,2))).evalf(30)-Q(626,100)) < Q(1,100) | so $e^{0.2t}=3.5$ and $t=5\ln3.5\approx6.26$. The domain of the solution is $[0,5\ln3.5]$: after that the formula gives negative amounts, which have no meaning.

\begin{center}
\begin{tikzpicture}[x=0.7cm,y=0.0045cm]
  \draw[->] (-0.2,0) -- (8.3,0) node[right] {$t$};
  \draw[->] (0,-80) -- (0,1200) node[above] {$p$};
  \draw[very thick,blue!70!black,domain=0:7,samples=40,variable=\t] plot ({\t},{200*exp(0.2*\t)+300});
  \draw[very thick,black!70,domain=0:7.5,samples=2,variable=\t] plot ({\t},{500});
  \draw[very thick,red!70!black,domain=0:6.2638,samples=40,variable=\t] plot ({\t},{700-200*exp(0.2*\t)});
  \fill (0,500) circle (1.6pt) node[left,font=\small] {$500$};
%@ g-e9-a | 200*exp(0)+300 == 500 and 700-200*exp(0) == 500 | circle (1.6pt) node[left,font=\small] {$500$};
  \fill (6.2638,0) circle (1.6pt) node[below,font=\small] {$6.26$};
%@ g-e9-b | 700-200*exp(5*log(Q(7,2))/5) == 0 | circle (1.6pt) node[below,font=\small] {$6.26$};
  \node[font=\small,blue!70!black,anchor=west] at (7,1110) {$H=60$};
  \node[font=\small,anchor=west] at (7.5,540) {$H=100$};
  \node[font=\small,red!70!black,anchor=west] at (5.3,270) {$H=140$};
\end{tikzpicture}
\par\small\emph{Figure.} $p'=0.2p-H$ with $p(0)=500$: growth for $H=60$, equilibrium for $H=100$, collapse at $t\approx6.26$ for $H=140$.
\end{center}
\end{exmp}

\begin{exmp}{Flow from a tank (Torricelli's law)}{e10}
A cylindrical tank of cross-section $A$ drains through a hole of area $a$ in the bottom. By Torricelli's law the depth $h(t)$ of water satisfies $h'(t)=-k\sqrt{h(t)}$, where $k=\frac aA\sqrt{2g}$ and $g=9.8$ m/s$^2$. (a) Verify that $h(t)=\left(\sqrt H-\frac{kt}2\right)^2$ solves the IVP with $h(0)=H$, and find its domain. (b) Take $H=0.64$ m, $A=0.8$ m$^2$ and $a=0.01$ m$^2$ and find when the tank is empty.

\textbf{Solution.} (a) Initial condition: $h(0)=(\sqrt H)^2=H$. Differentiate with the Chain Rule: $h'(t)=2\left(\sqrt H-\frac{kt}2\right)\left(-\frac k2\right)=-k\left(\sqrt H-\frac{kt}2\right)$. The equation requires $h'=-k\sqrt h$, and $\sqrt h=\sqrt{(\sqrt H-kt/2)^2}=\left|\sqrt H-\frac{kt}2\right|$. The two agree \emph{only when $\sqrt H-\frac{kt}2\ge0$}, that is for $t\le\dfrac{2\sqrt H}k$. For larger $t$ the formula is not a solution (its derivative is positive while $-k\sqrt h<0$; physically, water does not flow back in). So the domain is $\left[0,\dfrac{2\sqrt H}k\right]$.
(b) $\frac aA=\frac{0.01}{0.8}=\frac1{80}$ and $k=\frac{\sqrt{19.6}}{80}=\frac{7\sqrt{10}}{400}\approx0.0553$ m$^{1/2}$/s (because $\sqrt{19.6}=\sqrt{196/10}=14/\sqrt{10}$). With $\sqrt H=0.8$, $h(t)=\left(0.8-\frac{7\sqrt{10}}{800}t\right)^2$. The tank is empty at $T=\frac{2\sqrt H}k=\frac{1.6\cdot400}{7\sqrt{10}}=\frac{64\sqrt{10}}7\approx28.9$ s. The figure shows the part of the parabola that is the solution (solid) and the part that is not (dashed).
%@ e10-b | simplify(diff((Q(4,5)-7*sqrt(10)*t/800)**2,t) + (7*sqrt(10)/400)*(Q(4,5)-7*sqrt(10)*t/800)) == 0 and NUM(sqrt(Q(196,10))/80, 7*sqrt(10)/400) == 0 and NUM(2*Q(4,5)/(7*sqrt(10)/400), 64*sqrt(10)/7) == 0 and abs((64*sqrt(10)/7).evalf(30)-Q(2891,100)) < Q(1,100) | $T=\frac{2\sqrt H}k=\frac{1.6\cdot400}{7\sqrt{10}}=\frac{64\sqrt{10}}7\approx28.9$ s. The figure shows the part of the parabola that is the solution (solid) and the part that is not (dashed).

\begin{center}
\begin{tikzpicture}[x=0.15cm,y=7cm]
  \draw[->] (-1,0) -- (42,0) node[right] {$t$ (s)};
  \draw[->] (0,-0.03) -- (0,0.76) node[above] {$h$ (m)};
  \draw[very thick,blue!70!black,domain=0:28.9123,samples=40,variable=\t] plot ({\t},{(0.8-0.027670*\t)*(0.8-0.027670*\t)});
  \draw[thick,dashed,gray,domain=28.9123:40,samples=30,variable=\t] plot ({\t},{(0.8-0.027670*\t)*(0.8-0.027670*\t)});
  \fill (0,0.64) circle (1.6pt) node[left,font=\small] {$0.64$};
%@ g-e10-a | (Q(4,5)-7*sqrt(10)*0/800)**2 == Q(16,25) | circle (1.6pt) node[left,font=\small] {$0.64$};
  \fill (28.9123,0) circle (1.6pt) node[below,font=\small] {$28.9$};
%@ g-e10-b | simplify((Q(4,5)-7*sqrt(10)*(64*sqrt(10)/7)/800)**2) == 0 | circle (1.6pt) node[below,font=\small] {$28.9$};
\end{tikzpicture}
\par\small\emph{Figure.} The tank empties at $t\approx28.9$ s; the dashed part of the parabola is not a solution of $h'=-k\sqrt h$.
\end{center}
\end{exmp}

\begin{remark}[Existence and uniqueness]
Throughout this section, once we found a solution of an IVP we did not ask whether there could be another. Theorems~\ref{thm:anti} and~\ref{thm:nth} answer that for the equations $y^{(n)}=f(t)$: the solution exists and is unique on the whole interval where $f$ is continuous. For a general first-order IVP $y'=F(t,y)$, $y(a)=b$ the answer is an \emph{existence and uniqueness theorem}: if $F$ is ``well behaved'' near $(a,b)$, a solution exists and is unique on some interval around $a$. The proof is technical and belongs to a later course in differential equations; we state it here without proof and use it (as most textbooks do) for the IVPs of this chapter. Some hypothesis is needed: $y'=3y^{2/3}$ with $y(0)=0$ has the two solutions $y=0$ and $y=t^3$ (here $y^{2/3}$ means $\left(\sqrt[3]y\right)^2\ge0$; for $y=t^3$: $y'=3t^2$ and $3\left(t^3\right)^{2/3}=3t^2$), both on $(-\infty,\infty)$.
%@ rem-nu | all(diff(t**3,t).subs(t,v) == 3*(v**6)**Q(1,3) for v in (Q(1,2),Q(-3,1),Q(5,1))) | and $3\left(t^3\right)^{2/3}=3t^2$), both on $(-\infty,\infty)$.
\end{remark}

\begin{think}{Two solutions through one point}{thuniq}
The IVP $y'=3y^{2/3}$, $y(0)=0$ has both $y=0$ and $y=t^3$ as solutions. Does this contradict Theorem~\ref{thm:anti}(b)? Where does the right side $3y^{2/3}$ fail to be ``well behaved''?
\end{think}

% ---------------------------------------------------------------------
\subsection*{Example bank}

\begin{remark}[What's in this bank]
\begin{itemize}[nosep]
  \item \textbf{Group 1, Classifying and verifying:} order and linearity; verifying a family with Theorem~\ref{thm:lin}.
  \item \textbf{Group 2, General solutions by integration:} first- and third-order equations $y^{(n)}=f(t)$.
  \item \textbf{Group 3, Initial value problems:} a logarithmic right side; a second-order IVP that needs parts.
  \item \textbf{Group 4, Domains:} a rational right side (partial fractions) and a trigonometric one.
  \item \textbf{Group 5, Motion under gravity:} a stone thrown down a cliff; the same launch on two planets.
  \item \textbf{Group 6, Models checked by substitution:} cooling; logistic growth.
  \item \textbf{Group 7, Proofs:} a linearity proof; a proof that uses Theorem~\ref{thm:nth}.
\end{itemize}
\end{remark}

\subsection*{Group 1: Classifying and verifying}

\begin{exmp}{Order and linearity, with reasons}{b1}
State the order of each equation and whether it is linear, and give a reason. (a) $(t^2+1)y'-y=\sin t$. (b) $y''+t\,e^{y}=0$. (c) $y'''+5y''-\dfrac{y}{t}=t^3$. (d) $y'=\sqrt{y}$.

\textbf{Solution.} (a) Order $1$. The coefficients $t^2+1$ and $-1$ and the right side $\sin t$ depend only on $t$: \emph{linear}. (b) Order $2$. The unknown is inside the exponential, $e^y$: \emph{nonlinear}. (c) Order $3$. The terms $y''',\,y'',\,y/t$ are each (function of $t$)$\times$(derivative of $y$ to the first power): \emph{linear} (for $t\ne0$). (d) Order $1$. $\sqrt y=y^{1/2}$ is not a first power of $y$: \emph{nonlinear}.
\end{exmp}

\begin{exmp}{A family from a product}{b2}
Show that $y=e^{-t}\left(C_1\cos2t+C_2\sin2t\right)$ is a solution of $y''+2y'+5y=0$ for all constants.

\textbf{Solution.} The equation is linear, so by Theorem~\ref{thm:lin}(b) it is enough to check $y_1=e^{-t}\cos2t$ and $y_2=e^{-t}\sin2t$. Product rule: $y_1'=-e^{-t}\cos2t-2e^{-t}\sin2t$, and differentiating again, $y_1''=e^{-t}\cos2t+2e^{-t}\sin2t+2e^{-t}\sin2t-4e^{-t}\cos2t=-3e^{-t}\cos2t+4e^{-t}\sin2t$. Then
\[
y_1''+2y_1'+5y_1=e^{-t}\left[(-3-2+5)\cos2t+(4-4)\sin2t\right]=0 .
\]
For $y_2$: $y_2'=-e^{-t}\sin2t+2e^{-t}\cos2t$ and $y_2''=e^{-t}\sin2t-2e^{-t}\cos2t-2e^{-t}\cos2t-4e^{-t}\sin2t=-3e^{-t}\sin2t-4e^{-t}\cos2t$, so $y_2''+2y_2'+5y_2=e^{-t}\left[(-3-2+5)\sin2t+(-4+4)\cos2t\right]=0$.
%@ b2 | simplify(diff(exp(-t)*cos(2*t),t,2)+2*diff(exp(-t)*cos(2*t),t)+5*exp(-t)*cos(2*t)) == 0 and simplify(diff(exp(-t)*sin(2*t),t,2)+2*diff(exp(-t)*sin(2*t),t)+5*exp(-t)*sin(2*t)) == 0 and simplify(diff(exp(-t)*cos(2*t),t,2) - (-3*exp(-t)*cos(2*t)+4*exp(-t)*sin(2*t))) == 0 | so $y_2''+2y_2'+5y_2=e^{-t}\left[(-3-2+5)\sin2t+(-4+4)\cos2t\right]=0$.
Both vanish, so the whole family is a solution on $(-\infty,\infty)$.
\end{exmp}

\subsection*{Group 2: General solutions by integration}

\begin{exmp}{A first-order equation}{b3}
Find the general solution of $y'(x)=5x^4-\dfrac{6}{x^3}+3\sin\dfrac x2$ on $(0,\infty)$.

\textbf{Solution.} Integrate term by term (Theorem~\ref{thm:anti}): $\int5x^4\,dx=x^5$; $\int-6x^{-3}\,dx=3x^{-2}$; $\int3\sin\frac x2\,dx=-6\cos\frac x2$ (substitution $u=\frac x2$, Theorem~\ref{8.1-thm:subst}). So $y(x)=x^5+\dfrac3{x^2}-6\cos\dfrac x2+C$. Check: $y'=5x^4-6x^{-3}+3\sin\frac x2$. We use $(0,\infty)$ because $x^{-3}$ is undefined at $0$.
%@ b3 | ANT(x**5+3/x**2-6*cos(x/2), 5*x**4-6/x**3+3*sin(x/2), [Q(1,2),Q(3,1),Q(7,3)]) == 0 | So $y(x)=x^5+\dfrac3{x^2}-6\cos\dfrac x2+C$. Check: $y'=5x^4-6x^{-3}+3\sin\frac x2$. We use $(0,\infty)$ because $x^{-3}$ is undefined at $0$.
\end{exmp}

\begin{exmp}{A third-order equation}{b4}
Find the general solution of $y'''(t)=24t-6$ and the solution with $y(1)=2$, $y'(1)=-1$, $y''(1)=6$.

\textbf{Solution.} Integrate three times: $y''=12t^2-6t+C_1$; $y'=4t^3-3t^2+C_1t+C_2$; $y=t^4-t^3+\frac{C_1}2t^2+C_2t+C_3$. Three integrations, three constants (Theorem~\ref{thm:nth}). Conditions: $y''(1)=12-6+C_1=6$ gives $C_1=0$; $y'(1)=4-3+0+C_2=-1$ gives $C_2=-2$; $y(1)=1-1+0-2+C_3=2$ gives $C_3=4$. So $y(t)=t^4-t^3-2t+4$. Check: $y'=4t^3-3t^2-2$, $y''=12t^2-6t$, $y'''=24t-6$, and $y(1)=2$, $y'(1)=-1$, $y''(1)=6$.
%@ b4 | simplify(diff(t**4-t**3-2*t+4,t,3) - (24*t-6)) == 0 and (t**4-t**3-2*t+4).subs(t,1) == 2 and diff(t**4-t**3-2*t+4,t).subs(t,1) == -1 and diff(t**4-t**3-2*t+4,t,2).subs(t,1) == 6 | $y'''=24t-6$, and $y(1)=2$, $y'(1)=-1$, $y''(1)=6$.
\end{exmp}

\subsection*{Group 3: Initial value problems}

\begin{exmp}{A logarithmic solution}{b5}
Solve $y'(x)=\dfrac{6x}{x^2+1}$, $y(0)=2$.

\textbf{Solution.} With $u=x^2+1$, $du=2x\,dx$ (Theorem~\ref{8.1-thm:subst}): $\int\frac{6x}{x^2+1}\,dx=3\int\frac{du}u=3\ln u+C=3\ln(x^2+1)+C$ (no absolute value, since $x^2+1>0$). The condition: $y(0)=3\ln1+C=C=2$. So $y(x)=3\ln(x^2+1)+2$, on $(-\infty,\infty)$. Check: $y'=3\cdot\frac{2x}{x^2+1}=\frac{6x}{x^2+1}$.
%@ b5 | ANT(3*log(x**2+1)+2, 6*x/(x**2+1), [Q(1,2),Q(3,1),Q(-2,1)]) == 0 and 3*log(1)+2 == 2 | So $y(x)=3\ln(x^2+1)+2$, on $(-\infty,\infty)$. Check: $y'=3\cdot\frac{2x}{x^2+1}=\frac{6x}{x^2+1}$.
\end{exmp}

\begin{exmp}{A second-order problem with integration by parts}{b6}
Solve $y''(t)=te^{-t}$, $y(0)=1$, $y'(0)=0$.

\textbf{Solution.} \emph{First integration.} By parts (Theorem~\ref{8.2-thm:parts}) with $u=t$, $dv=e^{-t}dt$: $\int te^{-t}dt=-te^{-t}+\int e^{-t}dt=-te^{-t}-e^{-t}=-(t+1)e^{-t}$. So $y'=-(t+1)e^{-t}+C_1$, and $y'(0)=-1+C_1=0$ gives $C_1=1$: $y'=1-(t+1)e^{-t}$.
\emph{Second integration.} $\int(t+1)e^{-t}dt=-(t+1)e^{-t}+\int e^{-t}dt=-(t+2)e^{-t}$ (parts again, $u=t+1$). So $y=t+(t+2)e^{-t}+C_2$, and $y(0)=0+2+C_2=1$ gives $C_2=-1$:
\[
y(t)=t+(t+2)e^{-t}-1 .
\]
Check: $y'=1+e^{-t}-(t+2)e^{-t}=1-(t+1)e^{-t}$ and $y''=-e^{-t}+(t+1)e^{-t}=te^{-t}$; $y(0)=2-1=1$ and $y'(0)=1-1=0$.
%@ b6 | simplify(diff(t+(t+2)*exp(-t)-1,t,2) - t*exp(-t)) == 0 and (t+(t+2)*exp(-t)-1).subs(t,0) == 1 and diff(t+(t+2)*exp(-t)-1,t).subs(t,0) == 0 | and $y''=-e^{-t}+(t+1)e^{-t}=te^{-t}$; $y(0)=2-1=1$ and $y'(0)=1-1=0$.
\end{exmp}

\subsection*{Group 4: Domains}

\begin{exmp}{A rational right side}{b7}
Solve $y'(x)=\dfrac1{x^2-4}$, $y(3)=0$, and find the domain.

\textbf{Solution.} By partial fractions (Section~8.5), $\dfrac1{x^2-4}=\dfrac14\left(\dfrac1{x-2}-\dfrac1{x+2}\right)$, because $\frac14\cdot\frac{(x+2)-(x-2)}{x^2-4}=\frac14\cdot\frac4{x^2-4}$. Hence $y=\frac14\ln\abs{x-2}-\frac14\ln\abs{x+2}+C=\frac14\ln\left|\frac{x-2}{x+2}\right|+C$. The condition: $y(3)=\frac14\ln\frac15+C=0$, so $C=\frac14\ln5$. For $x>2$ the quotient is positive, so
\[
y(x)=\frac14\ln\frac{5(x-2)}{x+2} .
\]
The right side is undefined at $x=\pm2$, and $x_0=3$ lies in $(2,\infty)$, so the domain is $(2,\infty)$. Check: $y'=\frac14\left(\frac1{x-2}-\frac1{x+2}\right)=\frac1{x^2-4}$.
%@ b7 | ANT(log(5*(x-2)/(x+2))/4, 1/(x**2-4), [Q(3,1),Q(5,2),Q(7,1)]) == 0 and NUM(log(5*(3-2)/(3+2))/4, 0) == 0 and expand(Q(1,4)*((x+2)-(x-2))-1) == 0 | Check: $y'=\frac14\left(\frac1{x-2}-\frac1{x+2}\right)=\frac1{x^2-4}$.
\end{exmp}

\begin{exmp}{A secant derivative}{b8}
Solve $y'(t)=3\sec t\tan t$, $y(2\pi)=5$, and find the domain.

\textbf{Solution.} Since $\frac d{dt}\sec t=\sec t\tan t$, $y=3\sec t+C$. The condition: $y(2\pi)=3\sec2\pi+C=3+C=5$, so $C=2$ and $y(t)=3\sec t+2$. Check: $y'=3\sec t\tan t$. Now $\sec t$ is undefined where $\cos t=0$, at $t=\frac{3\pi}2$ and $\frac{5\pi}2$ among others; $2\pi$ lies between them, so the domain is $\left(\frac{3\pi}2,\frac{5\pi}2\right)$.
%@ b8 | ANT(3*sec(x)+2, 3*sec(x)*tan(x), [Q(1,2),Q(3,1),Q(5,1)]) == 0 and 3*sec(2*pi)+2 == 5 and cos(3*pi/2) == 0 and cos(5*pi/2) == 0 | so the domain is $\left(\frac{3\pi}2,\frac{5\pi}2\right)$.
\end{exmp}

\subsection*{Group 5: Motion under gravity}

\begin{exmp}{A stone thrown down a cliff}{b9}
A stone is thrown downward with initial speed $4.9$ m/s from the top of a $100$ m cliff. Find $s(t)$ (height above the ground), the time it hits the ground, and its speed then.

\textbf{Solution.} Upward is positive, so $v_0=-4.9$ and $s_0=100$. By Example~\ref{ex:e8}, $s(t)=-4.9t^2-4.9t+100$ and $v(t)=-9.8t-4.9$. Impact: $s=0$, so $4.9t^2+4.9t-100=0$, that is $t^2+t-\frac{1000}{49}=0$. By the quadratic formula the positive root is $t=\dfrac{-1+\sqrt{1+4000/49}}2=\dfrac{-7+\sqrt{4049}}{14}\approx4.05$ s. The velocity then is $v\approx-9.8(4.045)-4.9\approx-44.5$ m/s: the speed is about $44.5$ m/s.
%@ b9 | simplify(-Q(49,10)*((-7+sqrt(4049))/14)**2-Q(49,10)*(-7+sqrt(4049))/14+100) == 0 and abs(((-Q(98,10))*((-7+sqrt(4049))/14)-Q(49,10)).evalf(30)+Q(4454,100)) < Q(1,100) | The velocity then is $v\approx-9.8(4.045)-4.9\approx-44.5$ m/s: the speed is about $44.5$ m/s.
\end{exmp}

\begin{exmp}{The same launch on two planets}{b10}
A rover throws a stone upward at $14.7$ m/s from the ground on Earth ($g=9.8$) and on Mars ($g=3.7$ m/s$^2$). Compare the times and heights of the highest points and the landing times.

\textbf{Solution.} With $s''=-g$, $s(0)=0$, $s'(0)=v_0$ the solution is $s=-\frac g2t^2+v_0t$ (as in Example~\ref{ex:e8} with $g$ in place of $9.8$). The peak is at $t=\frac{v_0}g$ with height $\frac{v_0^2}{2g}$; the landing time is $\frac{2v_0}g$ ($s=t\left(v_0-\frac g2t\right)=0$). Earth: peak at $\frac{14.7}{9.8}=1.5$ s, height $\frac{14.7^2}{19.6}=11.025$ m, landing at $3$ s. Mars: peak at $\frac{14.7}{3.7}\approx3.97$ s, height $\frac{14.7^2}{7.4}\approx29.2$ m, landing at about $7.95$ s. The height is inversely proportional to $g$: Mars' $g$ is smaller by the factor $\frac{3.7}{9.8}$, so the stone goes $\frac{9.8}{3.7}\approx2.65$ times higher.
%@ b10 | Q(147,10)/Q(98,10) == Q(3,2) and Q(147,10)**2/Q(196,10) == Q(21609,1960) and Q(21609,1960) == Q(11025,1000) and abs((Q(147,10)**2/Q(74,10)).evalf(30)-Q(292,10)) < Q(1,10) and abs((Q(98,10)/Q(37,10)).evalf(30)-Q(265,100)) < Q(1,100) | so the stone goes $\frac{9.8}{3.7}\approx2.65$ times higher.
\end{exmp}

\subsection*{Group 6: Models checked by substitution}

\begin{exmp}{Newton's law of cooling}{b11}
The temperature $T(t)$ of a cup of coffee in a room at $A=22^\circ$C obeys $T'=-k(T-A)$. (a) Verify that $T(t)=(T_0-A)e^{-kt}+A$ solves it with $T(0)=T_0$. (b) With $T_0=90$ and $k=0.1$ per minute, when does the coffee reach $40^\circ$C?

\textbf{Solution.} (a) $T'(t)=-k(T_0-A)e^{-kt}$ and $-k(T-A)=-k(T_0-A)e^{-kt}$: equal. And $T(0)=T_0-A+A=T_0$. (b) $T(t)=68e^{-0.1t}+22=40$ gives $e^{-0.1t}=\frac{18}{68}=\frac9{34}$, so $t=10\ln\frac{34}9\approx13.3$ minutes. As $t\to\infty$, $T\to22$.
%@ b11 | simplify(diff((a-b)*exp(-k*t)+b,t) + k*((a-b)*exp(-k*t)+b-b)) == 0 and Q(18,68) == Q(9,34) and abs((10*log(Q(34,9))).evalf(30)-Q(133,10)) < Q(1,10) and 68*exp(-Q(1,10)*10*log(Q(34,9)))+22 == 40 | so $t=10\ln\frac{34}9\approx13.3$ minutes. As $t\to\infty$, $T\to22$.
\end{exmp}

\begin{exmp}{Logistic growth}{b12}
Show that $P(t)=\dfrac{400}{1+7e^{-0.3t}}$ solves $P'=0.3P\left(1-\dfrac P{400}\right)$, and find $P(0)$, the time at which $P=200$, and the largest value of $P'$.

\textbf{Solution.} Let $E=e^{-0.3t}$, so $E'=-0.3E$. Then $P=400(1+7E)^{-1}$ and $P'=-400(1+7E)^{-2}\cdot7E'=\dfrac{840E}{(1+7E)^2}$. On the other side, $1-\frac P{400}=1-\frac1{1+7E}=\frac{7E}{1+7E}$, so $0.3P\left(1-\frac P{400}\right)=0.3\cdot\frac{400}{1+7E}\cdot\frac{7E}{1+7E}=\frac{840E}{(1+7E)^2}$. Equal. Next, $P(0)=\frac{400}{1+7}=50$. $P=200$ when $1+7E=2$, so $E=\frac17$ and $t=\frac{\ln7}{0.3}\approx6.49$. The growth rate $0.3P(1-P/400)$ is a downward parabola in $P$ with zeros at $0$ and $400$, so its largest value is at $P=200$: $0.3\cdot200\cdot\frac12=30$, reached at $t\approx6.49$.
%@ b12 | simplify(diff(400/(1+7*exp(-Q(3,10)*t)),t) - Q(3,10)*(400/(1+7*exp(-Q(3,10)*t)))*(1-(400/(1+7*exp(-Q(3,10)*t)))/400)) == 0 and 400/(1+7*exp(0)) == 50 and abs((log(7)/Q(3,10)).evalf(30)-Q(649,100)) < Q(1,100) and Q(3,10)*200*(1-Q(1,2)) == 30 | so its largest value is at $P=200$: $0.3\cdot200\cdot\frac12=30$, reached at $t\approx6.49$.
\end{exmp}

\subsection*{Group 7: Proofs}

\begin{exmp}{Linear combinations (unwind a definition)}{b13}
Prove: if $y_1,y_2$ both solve $y''+p(t)\,y'+q(t)\,y=0$ on an interval $I$, then so does $3y_1-2y_2$. Check with $y''+9y=0$, $y_1=\cos3t$, $y_2=\sin3t$.

\textbf{Plan.} Write the equation as $L[y]=0$ and use the sum and constant-multiple rules.
\textbf{Proof.} Let $y=3y_1-2y_2$. Then $y'=3y_1'-2y_2'$ and $y''=3y_1''-2y_2''$. So $y''+py'+qy=3\left(y_1''+py_1'+qy_1\right)-2\left(y_2''+py_2'+qy_2\right)=3\cdot0-2\cdot0=0$ on $I$.
\emph{Check.} $y=3\cos3t-2\sin3t$: $y''=-27\cos3t+18\sin3t=-9y$, so $y''+9y=0$. \emph{Common mistake:} the same argument does not work for $y''+y^2=0$, because $(3y_1-2y_2)^2\ne3y_1^2-2y_2^2$.
%@ b13 | simplify(diff(3*cos(3*t)-2*sin(3*t),t,2)+9*(3*cos(3*t)-2*sin(3*t))) == 0 and simplify(diff(3*cos(3*t)-2*sin(3*t),t,2) - (-27*cos(3*t)+18*sin(3*t))) == 0 | because $(3y_1-2y_2)^2\ne3y_1^2-2y_2^2$.
\end{exmp}

\begin{exmp}{Solutions of $y'''=0$ (use an earlier theorem)}{b14}
Prove that every solution of $y'''=0$ on $(-\infty,\infty)$ is a polynomial of degree at most $2$. Find the solution with $y(1)=2$, $y'(1)=-1$, $y''(1)=6$.

\textbf{Plan.} Show that for any prescribed values the polynomial family contains a solution, then use the uniqueness in Theorem~\ref{thm:nth}.
\textbf{Proof.} Let $y$ be any solution and put $b_0=y(1)$, $b_1=y'(1)$, $b_2=y''(1)$. The polynomial $P(t)=b_0+b_1(t-1)+\frac{b_2}2(t-1)^2$ satisfies $P'''=0$, $P(1)=b_0$, $P'(1)=b_1$, $P''(1)=b_2$. So $P$ and $y$ solve the same IVP with $f=0$ (continuous on $(-\infty,\infty)$), and by the uniqueness in Theorem~\ref{thm:nth}, $y=P$.
\emph{The IVP:} $P(t)=2-(t-1)+3(t-1)^2=3t^2-7t+6$. \emph{Check:} $P(1)=3-7+6=2$, $P'(t)=6t-7$ gives $P'(1)=-1$, $P''=6$.
%@ b14 | expand(2-(t-1)+3*(t-1)**2-(3*t**2-7*t+6)) == 0 and (3*t**2-7*t+6).subs(t,1) == 2 and diff(3*t**2-7*t+6,t).subs(t,1) == -1 and diff(3*t**2-7*t+6,t,2) == 6 and diff(3*t**2-7*t+6,t,3) == 0 | \emph{Check:} $P(1)=3-7+6=2$, $P'(t)=6t-7$ gives $P'(1)=-1$, $P''=6$.
\end{exmp}

% ---------------------------------------------------------------------
\subsection*{Answers and discussion}

\begin{description}

  \item[Knowledge Check \ref{kc:kcord}.] (i) Order $2$, linear. (ii) Order $1$, linear (the right side $y+t^2$ is first degree in $y$). (iii) Order $3$, nonlinear (the product $y\,y'$).

  \item[Knowledge Check \ref{kc:kciv}.] $y=-48e^{-t/4}+C$ and $y(0)=-48+C=-5$ give $C=43$: $y(t)=43-48e^{-t/4}$. As $t\to\infty$, $y\to43$ from below (since $-48e^{-t/4}<0$).
%@ ans-kc2 | ANT(43-48*exp(-x/4), 12*exp(-x/4), [Q(1,2),Q(3,1)]) == 0 and 43-48 == -5 and limit(43-48*exp(-x/4), x, oo) == 43 | $y\to43$ from below (since $-48e^{-t/4}<0$).

  \item[Knowledge Check \ref{kc:kcdom}.] $y=3t+3t^{2/3}+C$ and $y(-8)=-24+3\cdot4+C=-12+C=10$ give $C=22$: $y(t)=3t+3t^{2/3}+22$. The initial point $-8$ lies in $(-\infty,0)$, and $y'$ is undefined at $t=0$, so the domain is $(-\infty,0)$ (it is the left branch).
%@ ans-kc3 | ANT(3*x+3*(-x)**Q(2,3)+22, 3-2*(-x)**Q(-1,3), [Q(-1,2),Q(-8,1)]) == 0 and 3*(-8)+3*4+22 == 10 | so the domain is $(-\infty,0)$ (it is the left branch).

  \item[Pause and think \ref{think:thgap}.] Yes: on $t<0$ the derivative is $\frac1{t^2}$, and on $t>0$ it is $\frac1{t^2}$, so $y'=\frac1{t^2}$ at every $t\ne0$. The constants may differ because there is no connection between the two pieces: the proof of Theorem~\ref{thm:anti}(a) uses the Mean Value Theorem between two points of the domain, which needs the whole segment between them to lie in the domain. On $(-\infty,0)\cup(0,\infty)$ two points on opposite sides do not have that. So the theorem is stated on an interval; on a union of intervals, the constant is chosen separately on each.
%@ ans-th1 | diff(-1/t+4,t) == 1/t**2 and diff(-1/t-7,t) == 1/t**2 | So the theorem is stated on an interval; on a union of intervals, the constant is chosen separately on each.

  \item[Pause and think \ref{think:thnl}.] $2y_1=\frac2{1-t}$ has $(2y_1)'=\frac2{(1-t)^2}$ but $(2y_1)^2=\frac4{(1-t)^2}$: not equal, so $2y_1$ is not a solution. In the proof of Theorem~\ref{thm:lin}(a) the step ``multiply the identity by $p_j$ and add'' works because $L$ is built from $y$ and its derivatives to the first power. For $y'=y^2$ the right side is $(cy_1)^2=c^2y_1^2$, not $c\,y_1^2$: the constant multiple does not come out unchanged, and (a) fails.
%@ ans-th2 | simplify(diff(2/(1-t),t) - 2/(1-t)**2) == 0 and simplify((2/(1-t))**2 - 4/(1-t)**2) == 0 | not equal, so $2y_1$ is not a solution. In the proof of Theorem~\ref{thm:lin}(a) the step

  \item[Pause and think \ref{think:thuniq}.] No. Theorem~\ref{thm:anti}(b) is about equations $y'=f(t)$ whose right side does not depend on $y$; here $F(t,y)=3y^{2/3}$ depends on $y$. The existence and uniqueness theorem asks for a ``well-behaved'' $F$. The function $y^{2/3}$ has the derivative $\frac23y^{-1/3}$, which is infinite at $y=0$: the slope of the direction field changes too violently there, and a solution can leave $y=0$ at any moment (for example $y=0$ for $t\le c$ and $y=(t-c)^3$ for $t>c$).
%@ ans-th3 | all(diff((t-1)**3,t).subs(t,v) == 3*((v-1)**6)**Q(1,3) for v in (Q(1,2),Q(-3,1),Q(5,1))) | at any moment (for example $y=0$ for $t\le c$ and $y=(t-c)^3$ for $t>c$).

\end{description}

\end{document}
